% Verbatim active mathematical statements from ../STOC2027.tex.
% Source SHA-256: b446db65a769b69d0968eaef6d9366ba94e14faa6c3fa7769348bd1f4c53b23a
% Reference appendix only; not a standalone LaTeX document.
% Commented, iffalsed, and junk-wrapped statements are excluded.

% STOC2027.tex:336-350
\begin{theorem}\label{thm:main}
Let
$n\geq2$, $0<\gamma\leq 1/4$, $0<\varepsilon\leq 1$, and
$0<\delta,\beta<1/10$.  There is an $(\varepsilon, \delta)$ edge-DP algorithm that, on every
nonnegative weighted $n$-vertex graph $G$, runs in time polynomial in
$n$, $1/\gamma$, $1/\varepsilon$, $\log(1/\delta)$, and $\log(1/\beta)$, and outputs a
nonnegative weighted graph $\widetilde{G}$ satisfying, with probability at least $1-\beta$, simultaneously
for every $S\subseteq V$,
\begin{align*}
 \left|w_G(S,V\setminus S)-w_{\widetilde{G}}(S,V\setminus S)\right|
 \leq \gamma w_G(S,V\setminus S)  + O\left(\frac{n}{\gamma\varepsilon}
 \left(\log(n/(\delta\beta\gamma\eps))\right)^{7}\right) .
\end{align*}
\end{theorem}

% STOC2027.tex:490-498
\begin{restatable}[Private correlation clustering]{corollary}{privatecorrelationclustering}
\label{cor:cc-intro}
For every constant $\varepsilon>0$ and $\delta=n^{-10}$, there is a
polynomial-time $(\varepsilon, \delta)$-edge-DP algorithm for unweighted complete-graph min-disagree
correlation clustering that, with high probability, outputs a partition with cost at most
\[
1.387\operatorname{OPT}_G+\widetilde O_{\varepsilon}(n).
\]
\end{restatable}

% STOC2027.tex:510-517
\begin{restatable}[Private Max-Cut]{corollary}{privatemaxcut}
\label{cor:maxcut-intro}
For every constant $\xi>0, \varepsilon>0$, and $\delta=n^{-10}$, there is a polynomial-time $(\varepsilon, \delta)$-edge-DP algorithm for Max-Cut that, for any non-negative weighted graph $G$, with high probability outputs a bipartition with cut value at least
\[
(0.878-\xi)\operatorname{OPT}_G-\widetilde O_{\xi,\varepsilon}(n).
\]
Here, $\operatorname{OPT}_G$ is the value of the real max-cut in $G$.
\end{restatable}

% STOC2027.tex:1458-1470
\begin{definition}
\label{def:DP}
Consider histograms $w,w' \in \R_{\geq0}^{\binom{n}{2}}$.
We say that histograms $w,w'$ are adjacent, written $w\sim w'$, if
$\norm{w-w'}_1\leq1$.
A randomized algorithm $\cM$ is $(\varepsilon, \delta)$ edge-DP if for
all adjacent $w,w'$ and measurable output events $A$,
\[
 \Prb[\cM(w)\in A]\leq e^\varepsilon\Prb[\cM(w')\in A]+\delta,
\]
and symmetrically with $w,w'$ exchanged.
\end{definition}

% STOC2027.tex:1490-1497
\begin{definition}[Laplace Distribution]
\label{def:laplace-dist}
The centered Laplace distribution with scale parameter $b > 0$, denoted as $\text{Lap}(b)$, is a continuous probability distribution with the probability density function
\[
    h(z) = \frac{1}{2b} \exp\left(-\frac{|z|}{b}\right), \quad z \in \mathbb{R}.
\]
If $Z \sim \text{Lap}(b)$, then $\Prb[|Z| \ge t] = \exp\left(-\frac{t}{b}\right)$ for $t\geq 0$.
\end{definition}

% STOC2027.tex:1521-1528
\begin{lemma}[Laplace Mechanism]
\label{lem:laplace-mech}
Let $f \colon \mathcal{X} \to \mathbb{R}^k$ be a function with $\ell_1$-sensitivity $\Delta_1 f = \max_{x \sim x'} \|f(x) - f(x')\|_1$. For any $\varepsilon > 0$, the Laplace mechanism $\mathcal{M}_L$ defined by
\[
    \mathcal{M}_L(x) = f(x) + (Z_1, \dots, Z_k),
\]
where $Z_1, \dots, Z_k$ are drawn i.i.d. from $\text{Lap}(\Delta_1 f / \varepsilon)$, satisfies $\varepsilon$-differential privacy.
\end{lemma}

% STOC2027.tex:1530-1533
\begin{lemma}[Post-Processing]
\label{lem:post-processing}
Let $\mathcal{M} \colon \mathcal{X} \to \mathcal{Y}$ be a randomized mechanism that is $(\varepsilon, \delta)$-differentially private. Let $f \colon \mathcal{Y} \to \mathcal{Z}$ be an arbitrary randomized mapping whose internal randomness is independent of the private input data. Then the composite mechanism $f \circ \mathcal{M} \colon \mathcal{X} \to \mathcal{Z}$ remains $(\varepsilon, \delta)$-differentially private.
\end{lemma}

% STOC2027.tex:1535-1538
\begin{lemma}[Basic Adaptive Composition]
\label{lem:basic-composition}
Let $\mathcal{M}_1, \dots, \mathcal{M}_k$ be a sequence of adaptively chosen randomized mechanisms, where for each $i \in [k]$, the mechanism $\mathcal{M}_i$ is $(\varepsilon_i, \delta_i)$-differentially private. Then their adaptive composition is $\left( \sum_{i=1}^k \varepsilon_i, \sum_{i=1}^k \delta_i \right)$-differentially private.
\end{lemma}

% STOC2027.tex:1540-1546
\begin{lemma}[Advanced Adaptive Composition~\cite{DRV10,DR14}]
\label{lem:advanced-composition}
Let $\mathcal{M}_1, \dots, \mathcal{M}_k$ be a sequence of adaptively chosen randomized mechanisms, where, conditional on every fixed preceding transcript, each $\mathcal{M}_i$ is $(\varepsilon_0, \delta_0)$-differentially private for some $\varepsilon_0 > 0$ and $\delta_0 \ge 0$. Then for any $\delta'\in(0,1)$, the adaptive composition of these $k$ mechanisms is $(\varepsilon, k\delta_0 + \delta')$-differentially private, where
$$
   \varepsilon = \varepsilon_0 \sqrt{2k \ln(1/\delta')} + k \varepsilon_0 (e^{\varepsilon_0} - 1).
$$
\end{lemma}

% STOC2027.tex:2003-2007
\begin{lemma}\label{lem:chi}
For every fixed $X \in \mathcal{X}$, choose independent random vectors $Z^{(j)}\sim \mathcal{N}(0,X)$ for each $j=1,\dots,k$. Set $\overline g(u,v)=\frac1k\sum_{j=1}^k\frac{(Z_u^{(j)}-Z_v^{(j)})^2}{4},$ and $\hat g(u,v)=\min\{\overline{g}(u,v)/2,1\}$. For some universal constant $C > 0$, if
$k\geq C\log(n/(\beta_\chi\del_1))$, then with probability at least $1-\min(\beta_\chi, \del_1)$, for every pair of nodes $e=\{u,v\}$ simultaneously,
\[\frac{q_{e}(X)}{4}\leq \hat{g}(u,v)\leq q_{e}(X).\]
\end{lemma}

% STOC2027.tex:2011-2017
\begin{lemma}\label{lem:chaos}
For every fixed $D \in \mathbb{R}^{\binom{n}{2}}$ and $X \in \mathcal{X}$, choose independent random vectors $Z^{(j)}\sim \mathcal{N}(0,X)$ for each $j=1,\dots,k$. Set $\overline g(u,v)=\frac1k\sum_{j=1}^k\frac{(Z_u^{(j)}-Z_v^{(j)})^2}{4},$ and $\hat g(u,v)=\min\{\overline{g}(u,v)/2,1\}$.  For some universal constant $C > 0$, if
$k\geq C\log(n/(\beta_\chi\del_1))$, for a universal constant $c_0>0$, conditioning on the good event in Lemma~\ref{lem:chi},
\[
 \Prb\left[\langle D,\hat{g}\rangle\geq\frac{1}{2}\sum_{e=\{u,v\} \in \binom{n}{2}} D_{uv} q_{e}(X)\right]\geq c_0.
\]
\end{lemma}

% STOC2027.tex:2189-2195
\begin{lemma}\label{lem:preprocess}
Fix any $\eps_0, \beta_0, \gamma \in (0, 1)$, and $n$-vertex input graph $G$ with edge-weight vector $w$. Decompose the weight vector $w=w'+v$, where $w'$ and $v$ are defined in \eqref{eq: def_split_weight}.
Then releasing $Y$ (as defined in \eqref{eq:noise_split_weight}) is $\eps_0$-DP, and, with probability at least $1-\beta_0$, for every cut $S$,
\[
 |Y(S,V\setminus S)+w'(S,V\setminus S)-w(S,V\setminus S)|\leq(\gamma/2)w(S,V\setminus S).
\]
\end{lemma}

% STOC2027.tex:2236-2246
\begin{lemma}\label{lem:stability_of_optimizer}
Let $A_1=A_0+\sum_e\xi_e b_eb_e^\top/4$, and let
\[
 X_i=\arg\max_{X\in\mathcal X}\{\ip{A_i}{X}+\lambda\log\det X\},
 \quad i\in\{0,1\}.
\]
Then
\[\norm{X_0^{-1/2}(X_1-X_0)X_1^{-1/2}}_F\leq \frac{\sum_e|\xi_e|\sqrt{q_e(X_0)q_e(X_1)}}{\lambda}.\]
Additionally, if $\norm{X_0^{-1/2}(X_1-X_0)X_1^{-1/2}}_F\leq1/2$, then for every $e$ it holds that
\[\frac{1}{2}q_e(X_1)\leq q_e(X_0)\leq 2q_e(X_1).\]
\end{lemma}

% STOC2027.tex:2309-2311
\begin{lemma}\label{lem:gaussian-privacy}
Let $\del_1 \in(0,1/2)$ be a fixed parameter and $X_0,X_1\in\mathbb{R}^{n\times n}$ be symmetric positive definite matrices. Denote $E=X_0^{-1/2}(X_1-X_0)X_1^{-1/2}$. We use $\pdf_0$ and $\pdf_1$ to denote the probability density functions of $\mathcal{N}(0,X_0)$ and $\mathcal{N}(0,X_1)$ respectively. Suppose $\norm{E}_F<1/2$. Then for \[\eps_1=O\left(k\norm{E}_F^2+\sqrt{k\log\frac{1}{\del_1}}\norm{E}_F+\log\frac{1}{\del_1}\norm{E}_F\right),\] we have \[e^{-\eps_1}\prod_{j=1}^k\pdf_1(Z^{(j)})\leq\prod_{j=1}^k\pdf_0(Z^{(j)})\leq e^{\eps_1}\prod_{j=1}^k\pdf_1(Z^{(j)})\] with probability at least $1-\del_1$ over the joint distribution of the batch   $(Z^{(1)}, \dots, Z^{(k)})$, where each $Z^{(j)} \sim \mathcal{N}(0, X_0)$ is drawn independently.
\end{lemma}

% STOC2027.tex:2317-2323
\begin{lemma}[\cite{laurent2000adaptive}]\label{lem:concentration_chi-squares}
For a fixed parameter $t>0$, given independent standard Gaussian variables $\Sigma_i\sim\mathcal{N}(0,1)$ and weight $a_i$ for $i=1,\dots,n$.  Then,
\[
 \Prb\left[\left|\sum_{i=1}^n a_i(\Sigma_i^2-1)\right|
 >2\bigl(\norm a_2\sqrt t+\norm a_\infty t\bigr)\right]\leq2e^{-t}.
\]
\end{lemma}

% STOC2027.tex:2391-2419
\begin{lemma}[Success-Charged Composition]
\label{lem:pay-what-you-get}
Suppose that the following conditions hold for every pair of adjacent inputs $w$ and $w'$. Let $\mathcal{M}_1, \dots, \mathcal{M}_k$ be a sequence of adaptively chosen randomized mechanisms, where each mechanism outputs either a valid result or a designated failure symbol $\bottom$.
For each step $t \in [k]$ and every common transcript prefix $y_{<t} = (y_1, \dots, y_{t-1})$, let $P_t^0$ and $P_t^1$ denote the conditional output distributions of $\mathcal{M}_t(w)$ and $\mathcal{M}_t(w')$, respectively. Suppose there exist auxiliary probability measures $N_t^0$ and $N_t^1$, a privacy loss parameter $\ell_t = \ell_t(y_{<t}) \ge 0$, and global constants $\ell^* < 1/10$ and $\delta' \ge 0$ such that, for any measurable set $A$ of return values and $b \in \{0, 1\}$:
\begin{itemize}
    \item If $P_t^b(A)=0$, then $N_t^b(A)=0;$
    \item $|P_t^b(A) - N_t^b(A)| \le \delta';$
    \item $e^{-\ell_t} N_t^1(A) \le N_t^0(A) \le e^{\ell_t} N_t^1(A),$
\end{itemize}
where $\ell_t(y_{<t}) \le \ell^*$.
Furthermore, suppose there exists a deterministic bound  $L$
such that for any sequence of outputs $y_1, \dots, y_k$ sampled sequentially according to
$y_t\sim N_t^b(\cdot\mid y_{<t}),$ for $t=1,\ldots,k$,
the accumulated privacy loss over successful attempts is bounded by $L$:
\begin{equation*}
    \sum_{t=1}^k \ell_t(y_{<t}) \cdot \ind[y_t \ne \bottom] \le L.
\end{equation*}
Then, for any fixed $\delta'' \in (0,1)$, the adaptive composition $(\mathcal{M}_1, \dots, \mathcal{M}_k)$ is $(\eps, \del)$-differentially private, where
\begin{equation*}
    \eps = 4L + 4\ell^*\log(1/\del''), \qquad \del = \del'' + (1+e^{\eps})k\del'.
\end{equation*}
\end{lemma}

% STOC2027.tex:2426-2431
\begin{lemma}[Privacy of \Cref{alg:gradient_aggregation}]
\label{lem:gradient_privacy}
Fix any $\eps_{\mathrm{GE}}, \delta_{\mathrm{GE}}$ and $\beta_{\mathrm{GE}}\in(0,1)$, let $\eps_1=\frac{\eps_{\mathrm{GE}}}{320K(1+\log(2/\del_{\mathrm{GE}}))}, \del_1=\frac{\del_{\mathrm{GE}}}{20n^2K}, k=\Theta(\log\frac{n}{\beta_{\chi}\del_1}), \beta_{\chi}=\beta_{\mathrm{GE}}/4n^2K$ and $\lambda=\Theta(\frac{1}{\eps_1}\log\frac{n}{\beta_{\mathrm{GE}}\del_{\mathrm{GE}}}).$ Then each invocation of
\Cref{alg:gradient_aggregation} in \Cref{alg:main} with input $(w_i,h_i,\gamma)$ is
$(\eps_{\mathrm{GE}}, \delta_{\mathrm{GE}})$-differentially private.
\end{lemma}

% STOC2027.tex:2655-2659
\begin{lemma}[Privacy of \Cref{alg:main}]\label{lem:main-privacy}
For $0<\varepsilon\leq1$ and $0<\delta<1$, fix $\eps_2=\frac{\eps}{2T_1}$, $\eps_{\mathrm{GE}}=\frac{\eps}{4T_1\sqrt{2T_2\log(2T_1/\delta)}}$, $\delta_{\mathrm{GE}}=\frac{\delta}{2T_1T_2}.$
Then \Cref{alg:main} is
$(\eps, \del)$-differentially private.
\end{lemma}

% STOC2027.tex:2699-2704
\begin{lemma}\label{lem:sdp_error}
For every signed vector $D \in \mathbb{R}^{\binom{n}{2}}$ and every cut $S \subseteq [n]$,
\[\langle D, c_S\rangle\leq 2\max\{F(D),0\}+O(\lambda n\log n+\norm{D}_1/n).\]
Moreover, $\langle D,\one\rangle\leq2F(D)$.
\end{lemma}

% STOC2027.tex:2722-2733
\begin{lemma}\label{lem:gradient_estimation}
Let $0<\gamma<1$, $K=\Theta(\log(n/\beta_{\mathrm{GE}}))$, $\tau_1=\log(1/\beta_1)/\eps_1$ and $\beta_1=\beta_{\chi}=\beta_{\mathrm{GE}}/4n^2K$. Given nonnegative edge weights $w$ and a nonnegative partial estimate $h$, suppose $w_e\leq U$ for every $e$. Let $(g^+,g^-,r^*)=\mathsf{AggregateWitnesses}(w,h,\gamma)$ (\Cref{alg:gradient_aggregation}).
Then with probability at least $1-\beta_{\mathrm{GE}}$:
\begin{itemize}
    \item $\langle h,g^+\rangle\geq(1+\gamma)\langle w,g^+\rangle;$
    \item $\langle h,g^-\rangle\leq(1-\gamma)\langle w,g^-\rangle;$
    \item For every cut $S$ simultaneously, \[(1-\gamma)\langle w,c_S\odotc r^*\rangle-B\leq \langle h,c_S\odotc r^*\rangle\leq(1+\gamma)\langle w,c_S\odotc r^*\rangle+B,\]
where $B=\Theta\left(\tau_1+\lambda n\log n+nU\right).$
\end{itemize}
\end{lemma}

% STOC2027.tex:2801-2809
\begin{lemma}\label{lem:inner_round}
Let $0 < \gamma \leq 1/4$, $\eta=\gamma/[8(1+\gamma)]$, $\tau_2=\log(1/\beta_2)/\eps_2$,
$B=O(\tau_1+\lambda n\log n+nU)$ as in \Cref{alg:main}, and
$M_{\mathrm{stop}}=C(nU\log\frac{nU}{\delta\beta\gamma\eps}+\frac{n}{\gamma\varepsilon}(\log\frac{nU}{\delta\beta\gamma\eps})^6)\geq10(B+\tau_2/\gamma)$, for a sufficiently large universal constant $C$.
For any fixed outer iteration $i$, with probability at least $1-T_2\beta_{\mathrm{GE}}-\beta_2$, if $\widetilde M_i>M_{\rm stop}$, then within at most $T_2=1+\left\lceil64(1+\gamma)^2\log n/\gamma^2\right\rceil=O(\log n/\gamma^2)$ inner iterations, $\mathsf{AggregateWitnesses}(w_i,h_i,\gamma)$ returns $(g^+_i,g^-_i,r_i)$ satisfying
\[
\langle h_i,r_i\rangle\geq\max\{\langle h_i,g^+_i\rangle,\langle h_i,g^-_i\rangle\}.
\]
\end{lemma}

% STOC2027.tex:2929-2931
\begin{lemma}\label{lem:outer_round}
With probability at least $1-T_1T_2\beta_{\mathrm{GE}}-T_1\beta_2$, \Cref{alg:main} terminates with $a_i=0$ or $\widetilde{M}_i\leq M_{\mathsf{stop}}$ in at most $T_1 = O(\log(nU))$ outer rounds.
\end{lemma}

% STOC2027.tex:2967-2983
\begin{lemma}\label{lem:utility}
Let $n\geq2$, $U\geq1$, $0<\gamma\leq1/4$, $0<\eps\leq1$, and $0<\delta,\beta<1/10$.
Choose all parameters as in \Cref{alg:main}, in particular
$\eps_{\mathrm{GE}}=\eps/(4T_1\sqrt{2T_2\log(2T_1/\delta)})$,
$\delta_{\mathrm{GE}}=\delta/(2T_1T_2)$, and $M_{\mathrm{stop}}=C(nU\log\frac{nU}{\delta\beta\gamma\eps}+\frac{n}{\gamma\varepsilon}(\log\frac{nU}{\delta\beta\gamma\eps})^6)
\geq10(B+\tau_2/\gamma)$, for a sufficiently large universal constant $C$.
Given a nonnegative weighted $n$-vertex graph $G$ with edge weights bounded by $U$, \Cref{alg:main} outputs a
nonnegative weighted graph $\widetilde{G}$ satisfying,
with probability at least $1-\beta$, simultaneously
for every $S\subseteq V$,
\begin{align*}
 \left|w_G(S,V\setminus S)-w_{\widetilde{G}}(S,V\setminus S)\right|
 \leq \gamma w_G(S,V\setminus S)  + O\left(nU\log\frac{nU}{\delta\beta\gamma\eps}+\frac{n}{\gamma\varepsilon}
 \left(\log\frac{nU}{\delta\beta\gamma\eps}\right)^7\right).
\end{align*}
\end{lemma}

% STOC2027.tex:3363-3368
\begin{lemma}[Chi-squared Tail Bound~\cite{laurent2000adaptive}]\label{lem:chi_square_tail}
Let $Y_1, \dots, Y_k$ be independent standard normal random variables, i.e., $Y_i \sim \mathcal{N}(0, 1)$. The random variable $Q = \sum_{i=1}^k Y_i^2$ follows a Chi-squared distribution with $k$ degrees of freedom, denoted as $Q \sim \chi_k^2$. For any such $Q$, there exists a universal constant $c > 0$ such that
\[
 \Prb\left[\frac{Q}{k} \notin \left[\frac{1}{2}, 2\right]\right] \leq 2e^{-ck}.
\]
\end{lemma}

% STOC2027.tex:3372-3378
\begin{lemma}[Gaussian Hypercontractivity for Degree-Two Polynomials~\cite{Janson97}]\label{lem:hypercontractivity}
Let $W$ be a degree-two polynomial in independent standard Gaussian random variables. Then its higher moments can be bounded by its second moment. Specifically, for the $4$-th moment, we have:
\[
    \|W\|_4 \le 3\|W\|_2,
\]
where $\|W\|_p = (\E[|W|^p])^{1/p}$ denotes the $L_p$ norm of the random variable.
\end{lemma}

% STOC2027.tex:3380-3386
\begin{lemma}[Log-Convexity of $L_p$ Norms]\label{lem:holder_interpolation}
For any random variable $W$ and $1 \le p < r < q \le \infty$, the $L_p$ norms satisfy the interpolation inequality $\|W\|_r \le \|W\|_p^\theta \|W\|_q^{1-\theta}$, where $\theta \in (0, 1)$ is determined by $\frac{1}{r} = \frac{\theta}{p} + \frac{1-\theta}{q}$.
In particular, choosing $p=1, r=2$, and $q=4$ yields $\theta = 1/3$, which gives:
\[
    \|W\|_2 \le \|W\|_1^{1/3} \|W\|_4^{2/3}.
\]
\end{lemma}

% STOC2027.tex:3390-3396
\begin{proposition}[Unbiased Estimator of Fractional Query]\label{prop:unbiased_estimator}
Let $X \in \mathcal{X}$ be a valid covariance matrix and let $Z \sim \mathcal{N}(0, X)$. For any pair of vertices $e = \{u,v\}$, we have
\[
    \E\left[\frac{(Z_u - Z_v)^2}{4}\right] = q_e(X).
\]
Consequently, for $\overline{g}(u,v) = \frac{1}{k} \sum_{j=1}^k \frac{(Z^{(j)}_u - Z^{(j)}_v)^2}{4}$ where $Z^{(j)}$ are drawn independently from $\mathcal{N}(0, X)$, it holds that $\E[\overline{g}(u,v)] = q_e(X)$.
\end{proposition}

% STOC2027.tex:3573-3586
\begin{lemma}[Generalized Pythagorean Inequality for KL Divergence~\cite{Csiszar75}]\label{lem:pythagorean}
Let $\mathcal{K}$ be a closed convex set of probability distributions, and let $\overline{p}$ be an arbitrary strictly positive probability distribution. If $p'$ is the projection of $\overline{p}$ onto $\mathcal{K}$, defined as:
\[
    p' = \argmin_{z \in \mathcal{K}} \KL(z \| \overline{p}),
\]
then for any target distribution $p^* \in \mathcal{K}$, it holds that:
\[
    \KL(p^* \| \overline{p}) \ge \KL(p^* \| p') + \KL(p' \| \overline{p}).
\]
In particular, since $\KL(p' \| \overline{p}) \ge 0$, this implies the distance-reduction property:
\[
    \KL(p^* \| p') \le \KL(p^* \| \overline{p}).
\]
\end{lemma}

% STOC2027.tex:3588-3599
\begin{lemma}[Capped unit-bounded cut release]\label{lem:capped-unit-release}
Let $D\subseteq\binom V2$ be public and let $x\in[0,1]^D$. For $0<\gamma\leq1/4$, $0<\varepsilon\leq1$, and $0<\delta,\beta_{\rm cap}<1/10$, there is a polynomial-time $(\varepsilon, \delta)$-DP mechanism under adjacency defined in Definition~\ref{def:DP} on $x$ that outputs $q\in[0,1]^D$ such that, except with probability $\beta_{\rm cap}$, simultaneously for every $S\subseteq V$,
\begin{equation*}
 \left|q(\delta_D(S))-x(\delta_D(S))\right| \leq \gamma x(\delta_D(S))+A_{\rm cap},
\end{equation*}
where
\[
 A_{\rm cap} \leq O\left(\frac{n}{\gamma\varepsilon}
 \left(\log\frac{n}{\delta\beta_{\rm cap}\gamma\varepsilon}\right)^7\right).
\]
\end{lemma}

% STOC2027.tex:3693-3703
\begin{theorem}[Black-box private correlation-clustering reduction]
\label{thm:cc-black-box-reduction}
Let $\rho\geq1$ and $p_0>0$ be constants, and let $\mathcal A$ be a polynomial-time randomized algorithm such that, on every unweighted complete-graph min-disagree instance $Q$,
\[
 \Prb\left[\operatorname{cc}_Q(\mathcal A(Q)) \leq \rho\operatorname{OPT}_Q\right] \geq p_0.
\]
Then, for every fixed $\xi>0$, $\varepsilon>0$, and $\delta=n^{-10}$, there is a polynomial-time $(\varepsilon, \delta)$-DP algorithm that, on every complete-graph min-disagree instance $G$, with high probability, outputs a partition $\widehat\Pi$ satisfying
\[
 \operatorname{cc}_G(\widehat\Pi) \leq (\rho+\xi)\operatorname{OPT}_G+\widetilde O_{\rho,\xi,\varepsilon}(n).
\]
\end{theorem}

% STOC2027.tex:3779-3790
\begin{corollary}[Private fixed-$k$ correlation clustering]
\label{cor:cc-fixed-k}
For every fixed $k\geq2$ and $\xi>0$, constant $\varepsilon>0$, and $\delta=n^{-10}$, there is a
polynomial-time $(\varepsilon, \delta)$-edge-DP algorithm for complete unweighted
Min-Disagree with at most $k$ clusters which, with high probability, outputs a partition
$\widehat\Pi$ satisfying
\[
 \operatorname{cc}_G(\widehat\Pi)
 \leq (1+\xi)\operatorname{OPT}_{G,k}
   +\widetilde O_{\xi,\varepsilon}(n).
\]
\end{corollary}

% STOC2027.tex:3824-3849
\begin{theorem}[Private Max-$k$-Cut]
\label{thm:private-max-k-cut}
Suppose a polynomial-time nonprivate algorithm, on every nonnegative weighted graph $Q$, outputs
an at-most-$k$ partition $\Pi$ and, for some $0<\rho_k\leq1$, satisfies
\[
 \E[\operatorname{cut}_Q(\Pi)]
 \geq\rho_k\operatorname{OPT}^{\operatorname{Max-}k\operatorname{-Cut}}_Q.
\]
For every fixed $\xi>0$, constant $\varepsilon>0$, and
$\delta=n^{-10}$, there is a polynomial-time $(\varepsilon, \delta)$-edge-DP algorithm which, with
high probability, outputs an at-most-$k$ partition $\widehat\Pi$ satisfying
\[
 \operatorname{cut}_G(\widehat\Pi)
 \geq (\rho_k-\xi)\operatorname{OPT}^{\operatorname{Max-}k\operatorname{-Cut}}_G
   -\widetilde O_{\xi,\varepsilon}(n).
\]
The additive term is independent of $k$.  In particular, for Max-Cut, the algorithm of
Goemans and Williamson~\cite{GoemansWilliamson95} gives
\[
 \operatorname{cut}_G(\widehat\Pi)
 \geq (0.878-\xi)\operatorname{OPT}^{\operatorname{Max-Cut}}_G
   -\widetilde O_{\xi,\varepsilon}(n).
\]
For $k\geq3$, one may use the standard SDP ratio $\rho_k$ of Frieze and
Jerrum~\cite{FriezeJerrum97}.
\end{theorem}

% STOC2027.tex:3893-3932
\begin{remark}[Approximate-DP Max-Cut lower bound]\label{rem:maxcut-tightness}
A direct-sum adaptation of the privacy comparison in~\cite{dharangutte2026price}
shows that the linear additive dependence on $n$ is necessary even for unweighted
Max-Cut under approximate differential privacy. Specifically, let $n\geq8$, $0\leq\varepsilon\leq1$,
$0\leq\delta,\beta\leq1/100$, and $0.868\leq\rho\leq1$.
If an $(\varepsilon,\delta)$-edge-DP algorithm $\mathcal M$ returns a cut of value
at least $\rho\operatorname{OPT}_G-B$, where $B\geq0$, with probability at least
$1-\beta$ on every unweighted $n$-vertex graph, then
$B\geq\lfloor n/8\rfloor/100=\Omega(n)$.
To see this, partition $8m$ vertices into $m=\lfloor n/8\rfloor$ public blocks
of size eight, and put one independently uniform random edge in each block;
any remaining vertices are isolated. Every such graph has optimum $m$.
Fix a block and the edges in all other blocks. Let $G_0$ leave this block empty,
and let $G_e=G_0+e$ for each of its $28$ possible edges $e$.
Let $E_e$ be the event that the output cut does not cut $e$.
Differential privacy gives
\[
 \Prb[\mathcal M(G_e)\in E_e]
 \geq e^{-\varepsilon}\bigl(\Prb[\mathcal M(G_0)\in E_e]-\delta\bigr).
\]
Every cut separates at most $16$ of the $28$ possible edges in the block, so
averaging over $e$ bounds its expected contribution under $\mathcal M(G_e)$ by
\[
 q:=1-e^{-\varepsilon}(3/7-\delta)
 \leq1-e^{-1}(3/7-1/100)<0.847.
\]
This holds for every choice of the other blocks. Averaging over all input edges
and summing over blocks therefore gives some graph $G$ with
$\E[\operatorname{cut}_G(\mathcal M(G))]\leq qm$.
The assumed utility guarantee and nonnegativity of cut values give
$\E[\operatorname{cut}_G(\mathcal M(G))]\geq(1-\beta)\rho m-B$.
Consequently,
\[
 B\geq\bigl((1-\beta)\rho-q\bigr)m
 >\bigl(0.99\cdot0.868-0.847\bigr)m>m/100.
\]
In particular, this applies to the ratio $0.878-\xi$ for $0<\xi\leq1/100$,
establishing the claimed tightness up to polylogarithmic factors in that range.
\end{remark}

% STOC2027.tex:3961-3991
\begin{theorem}[Private bicriteria cut minimization]
\label{thm:private-constrained-cut}
Fix $0<\eta\leq1/4$, $0<\varepsilon\leq1$, $0<\delta,\beta<1/10$,
$\rho\geq1$, and $0\leq\tau<1$.  Let
$\emptyset\ne\mathcal F,\mathcal F'\subseteq2^V$ be public families and
define
\[
 \operatorname{OPT}_{G,\mathcal F}
 :=\min_{S\in\mathcal F}w_G(S,V\setminus S).
\]
Suppose a polynomial-time nonprivate algorithm, on every nonnegative weighted graph, returns
$\widehat S\in\mathcal F'$ with probability at least $1-\tau$ and, on that event,
\[
 w(\widehat S,V\setminus\widehat S)
 \leq \rho\min_{S\in\mathcal F}w(S,V\setminus S).
\]
Then there is a polynomial-time $(\varepsilon, \delta)$-edge-DP algorithm which, with probability
at least $1-\beta-\tau$, returns $\widehat S\in\mathcal F'$ satisfying
\[
 w_G(\widehat S,V\setminus\widehat S)
 \leq \rho\frac{1+\eta}{1-\eta}\operatorname{OPT}_{G,\mathcal F}
   +\frac{\rho+1}{1-\eta}A_\eta.
\]
Here $A_\eta$ is the quantity in \eqref{eq:application-additive-error} for the displayed
$(\varepsilon, \delta,\beta)$.
Consequently, for every fixed $\xi>0$ and $\rho\geq1$, choosing
$\eta=\min\{1/4,\xi/(2\rho+\xi)\}$ gives
$(\rho+\xi)\operatorname{OPT}_{G,\mathcal F}
+\widetilde O_{\xi,\varepsilon}(\rho^{8}n)$ for constant $\varepsilon$ and
inverse-polynomial $\delta,\beta$.
\end{theorem}

% STOC2027.tex:4087-4099
\begin{lemma}[Bicriteria balanced recursion]\label{lem:hc-bicriteria-recursion}
Fix constants $0<b\leq1/3$ and $0<b'\leq1/3$, and let $a\geq1$.  Suppose an algorithm, on every
nonnegative weighted graph with vertex set $A$ of size at least $1/b$, returns a $b'$-balanced
cut $S$ satisfying
\[
 \phi_{G[A]}(S)\leq a\min\{\phi_{G[A]}(R):R\text{ is $b$-balanced}\}.
\]
Recursively applying this algorithm, and splitting smaller clusters as evenly as possible,
produces a hierarchical clustering tree $T$ with
\[
 \operatorname{hcost}_G(T)\leq O_b(a)\operatorname{OPT}^{\rm HC}_G.
\]
\end{lemma}

% STOC2027.tex:4150-4163
\begin{theorem}[Private hierarchical clustering]
\label{thm:private-hierarchical-clustering}
For every $0<\eta\leq1/4$, constant $0<\varepsilon\leq1$, and $\delta=n^{-10}$, there is a
polynomial-time $(\varepsilon, \delta)$-edge-DP algorithm that, with high probability, outputs a
$b'$-balanced hierarchical clustering tree $\widehat T$, for a universal constant $b'>0$,
satisfying
\[
 \operatorname{hcost}_G(\widehat T)
 \leq O(\sqrt{\log n})\operatorname{OPT}^{\rm HC}_G
  +O_\eta(A_\eta n\log^{3/2}n).
\]
In particular, for fixed $\eta$ the additive term is
$\widetilde O_{\eta,\varepsilon}(n^2)$.
\end{theorem}

% STOC2027.tex:4356-4367
\begin{lemma}[{\cite[Lemma~5.2]{EKKL20}}]
\label{lem:ekkl-discrepancy}
There is an absolute constant $0<\kappa\leq1$ such that the following holds.
Let $n\geq2$, $1\leq d\leq n/2$, and $\sigma\in[0,1]$.
Let $G,G'$ be simple unweighted graphs on the same $n$ vertices, with every
vertex having degree in $[d/2,2d]$ in each graph. If
$|E(G)\mathbin{\triangle}E(G')|\geq\sigma nd$, then
\[
 \max_{\substack{S,T\subseteq V\\S\cap T=\emptyset}}
       |w_G(S,T)-w_{G'}(S,T)|\geq\kappa\sigma n\sqrt d.
\]
\end{lemma}

% STOC2027.tex:4377-4397
\begin{theorem}[Lower Bound for Relative Cut Release]
\label{thm:lower_bound}
Let $n\geq33$, $\varepsilon\in(0,1)$, $0\leq\delta\leq\varepsilon/100$,
$\beta\in (0,1/40]$, and $\gamma \in [0, 1)$.
Suppose that an $(\varepsilon,\delta)$-differentially private mechanism
$\mathcal M$, on every nonnegative weighted graph $G$ on a fixed set $V$
of $n$ vertices, outputs a nonnegative weighted graph $H$ such that,
with probability at least $1-\beta$, simultaneously for all $S\subseteq V$,
\[
 \left|w_G(S, V\setminus S) - w_H(S, V\setminus S)\right| \leq \gamma w_G(S, V\setminus S) + A.
\]
Then the additive error must satisfy
\[
 A=\Omega\left(
   (1-\gamma)\cdot\min\left\{\frac{n^{3/2}}{\varepsilon},
                    \frac{n}{\varepsilon\gamma}\right\}
 \right).
\]
where the implicit constant is absolute and $1/\gamma$ is interpreted as
$+\infty$ when $\gamma=0$.
\end{theorem}
